\section{Artifact Access Semantics}
\label{sec:aas}

The \emph{Experience Compression Spectrum} (ECS) studies how an agent's interaction trace is transformed into reusable artifacts \cite{zhang2026experiencecompressionspectrumunifying}.
Following ECS, let an interaction trace be
\begin{equation}
T=\{(s_t,a_t,o_t,f_t)\}_{t=1}^{N},
\end{equation}
and let a level-$L$ compression function be
\begin{equation}
C_L:T\rightarrow K_L .
\end{equation}
Different choices of $L$ produce artifacts at different abstraction levels, such as raw traces, episodic memories, procedural skills, declarative rules, state summaries, or workspaces.
This axis answers a storage-side question: what does past interaction become?

We study the complementary runtime question: once past interaction has been compressed into artifacts, how does a current intent use them for a future decision?
A compressed artifact is not only a content payload.
It also carries usage semantics: a trace can be searched or replayed, a state can be queried, a graph can be traversed, a skill can be invoked, a rule can be injected or checked, and a workspace can be read, edited, or verified.
We call this runtime problem \emph{Artifact Access Semantics} (AAS).
Its core chain is
\begin{equation}
\begin{aligned}
\text{compression}&\rightarrow\text{artifact}
\rightarrow\text{usage semantics}\\
&\rightarrow\text{access}\rightarrow\text{decision}.
\end{aligned}
\end{equation}

\subsection{From Access Interfaces to Decision Views}
\label{subsec:aas-map}

AAS treats retrieval as one access interface rather than the general form of memory use.
Let $\mathcal{P}$ be the set of typed access primitives exposed by a memory system.
At runtime, a current intent $u$ induces an access trace
\begin{equation}
\tau_\pi(u,\mathcal{A})
=((\rho_1,\theta_1,r_1),\ldots,(\rho_T,\theta_T,r_T)),
\end{equation}
where $\rho_t$ is an access primitive, $\theta_t$ its parameters, and $r_t$ its returned result.
The trace may be a one-shot retrieval call, or a multi-step process involving graph traversal, source reading, workspace interaction, execution, or verification.
The returned observations are then organized into a decision-facing view $v=\mathcal{V}(u,\tau_\pi)$ and passed to a downstream model.

This abstraction separates three questions that are often collapsed in end-to-end memory QA: whether the relevant artifact exists, whether a budget-feasible access trace exposes it, and whether the exposed view is sufficient for the current decision.
Table~\ref{tab:aas-system-map} sketches how common systems occupy this space.

\begin{table*}[t]
\centering
\footnotesize
\setlength{\tabcolsep}{3.5pt}
\begin{tabularx}{\textwidth}{
>{\raggedright\arraybackslash}p{0.18\textwidth}
>{\raggedright\arraybackslash}p{0.20\textwidth}
>{\raggedright\arraybackslash}p{0.22\textwidth}
>{\raggedright\arraybackslash}X}
\toprule
Line of work & Artifact substrate & Runtime access interface & Exposed decision view \\
\midrule
Long-context agents
& Raw interaction history
& Full-context reading
& Long prompt context; history is present but not necessarily localized or updated. \\
RAG / memory retrieval
& Text chunks or memory records
& Top-$k$ retrieval and reranking
& Ranked evidence passages. \\
GraphRAG / hierarchical memory
& Entity, relation, or category graphs
& Graph traversal, neighborhood expansion, or global selection
& Graph-selected evidence; structure guides access. \\
Episodic / workspace memory
& Actor-state-event workspaces
& Entity, event, or state lookup
& Structured episodic summaries or state-like views. \\
Memory operating systems
& Persistent user or task states
& State lookup, update, and personalization hooks
& Personalized memory context. \\
Direct corpus / code agents
& Raw corpora, files, repositories
& Grep, read, shell, execution, navigation
& Composed evidence or verified result. \\
\bottomrule
\end{tabularx}
\caption{AAS view of runtime artifact access. The map distinguishes artifact substrates, access interfaces, and the decision views exposed downstream.}
\label{tab:aas-system-map}
\end{table*}

\subsection{Diagnostic Scope}
\label{subsec:aas-scope}

The rest of the paper focuses on the boundary after relevant artifacts, or oracle evidence, are available.
This is intentionally narrower than the full memory lifecycle.
We do not attempt to solve how memories are written, updated, indexed, or retrieved in general.
Instead, we ask whether available evidence is exposed in a decision-facing form and whether the decoder can use it.

We distinguish two downstream failures.
A \emph{view-use failure} occurs when needed evidence exists and can be accessed, but the exposed view omits decision-critical structure such as bindings, temporal authority, source relations, aggregation fields, modality-specific information, freshness, or conflict status.
A \emph{decoder-use failure} occurs when the view exposes the required information, but the decoder still fails to produce an evidence-valid answer.
This diagnostic scope motivates the theory in Section~\ref{sec:framework}: access can make evidence reachable, but post-access evaluation must make it answerable.
